% !TeX root = Report.tex
\chapter{Использование CSSA для предобработки изображений}

В данной главе мы приводим теоретическое обоснование методов на основе CSSA,
описанных в предыдущей главе. Сначала рассматривается вариант
\texttt{ROW.ROW}: показывается, как строка изображения после DFT сводится к
комплексно-экспоненциальному временному ряду, и объясняется точное восстановление в этой модели.
Затем рассматривается вариант \texttt{COL.ROW}, для которого существенную роль играет
ранг траекторной матрицы, построенной по столбцам комплексной матрицы. Отдельно
мы обсуждаем эффекты, возникающие после обратного преобразования Фурье, а также
два способа их уменьшения: ESPRIT и пороговую фильтрацию восстановленного
изображения.


\section{Теория \texorpdfstring{\texttt{CSSA ROW.ROW}}{CSSA ROW.ROW}}

Приведем несколько утверждений про математическую модель алгоритма. Пусть
$p$-я строка изображения содержит один пиксель сигнала в столбце $j_p$ с
интенсивностью $c_p>0$. Тогда соответствующий одномерный ряд имеет вид
\[
\mathbf{x}^{(p)} = c_p \mathbf{e}_{j_p} + \boldsymbol{\varepsilon}^{(p)},
\]
где $\mathbf{e}_{j_p}$ --- единичный орт длины $N_{\mathrm{col}}$,
а $\boldsymbol{\varepsilon}^{(p)}$ --- шумовая компонента. Для сигнала
после дискретного преобразования Фурье получаем
\begin{equation}
\mathcal{F}\bigl(c_p\mathbf{e}_{j_p}\bigr)
=
c_p \boldsymbol{\varphi}_{j_p},
\qquad
\bigl(\boldsymbol{\varphi}_{j_p}\bigr)_k
=
\exp\!\left(
-\frac{2\pi\mathrm{i}}{N_{\mathrm{col}}}(j_p-1)(k-1)
\right),
\quad k=1,\dots,N_{\mathrm{col}}.
    \label{eq:one-pixel-dft-exponential}
\end{equation}
Таким образом, строка с одним пикселем сигнала после DFT задается
одной комплексной экспонентой. Именно это является теоретическим основанием
для сохранения одной компоненты CSSA в варианте \texttt{ROW.ROW}.

\begin{statement}
Пусть задан комплексный ряд
\[
z_k = c\mu^{k-1}, \qquad k=1,\dots,N,
\]
где $c \ne0$, $\mu\in\mathbb{C}$, $\mu\ne0$. Пусть $1<L<N$ и
$K=N-L+1$. Тогда траекторная матрица
\[
\mathcal{H}_L(z)
=
\{z_{p+\ell-1}\}_{p=1,\ell=1}^{L,K}
\]
имеет ранг $1$.
\end{statement}

\begin{proof}
Элемент траекторной матрицы равен
\[
z_{p+\ell-1}=c\mu^{p+\ell-2}
=
c\mu^{p-1}\mu^{\ell-1}.
\]
Следовательно,
\[
\mathcal{H}_L(z)=uv^{\top},
\qquad
u_p=c\mu^{p-1},\quad v_\ell=\mu^{\ell-1}.
\]
В силу того, что
$\rank(\mathcal{H}_L(z)) \le \min(\rank(u), \rank(v)) = 1$ и
$\mathcal{H}_L(z)\ne0$, получаем $\rank(\mathcal{H}_L(z))=1$.
\end{proof}

\begin{corollary}
Для любого положения пикселя $j\in\{1,\dots,N_{\mathrm{col}}\}$ чистая строка
$c\mathbf{e}_j$, $c\ne0$, точно восстанавливается методом
\texttt{CSSA ROW.ROW} при сохранении одной ненулевой компоненты.
\end{corollary}

\begin{proof}
По формуле \eqref{eq:one-pixel-dft-exponential} после DFT строка
$c\mathbf{e}_j$ переходит в комплексную экспоненту.
В предыдущем утверждении достаточно взять
\[
\mu=\exp\!\left(
-\frac{2\pi\mathrm{i}}{N_{\mathrm{col}}}(j-1)
\right).
\]
Значит, траекторная матрица этого ряда имеет ранг $1$. Поэтому
Реконструкция после CSSA по единственной ненулевой компоненте восстанавливает ряд
$z$ точно, а обратное DFT возвращает исходную строку $c\mathbf{e}_j$.
\end{proof}

\begin{remark}
Если в строке находится несколько пикселей сигнала,
\[
\mathbf{x}^{(p)}
=
\sum_{\ell=1}^{q_p} c_{p\ell}\mathbf{e}_{j_{p\ell}},
\]
то, в силу линейности DFT, компонента сигнала становится суммой $q_p$
комплексных экспонент:
\[
\mathcal{F}\bigl(\mathbf{x}^{(p)}\bigr)
=
\sum_{\ell=1}^{q_p} c_{p\ell}\boldsymbol{\varphi}_{j_{p\ell}}.
\]
Поэтому ранг траекторной матрицы такого ряда не превосходит $q_p$.
\end{remark}

Итак, после DFT строка с одним пикселем сигнала имеет вид
\[
F\mathbf{x}^{(p)}
=
c_p\boldsymbol{\varphi}_{j_p}
+
\widetilde{\boldsymbol{\varepsilon}}^{(p)}.
\]
Компонента сигнала имеет ранг $1$, а шумовая компонента, вообще говоря, не имеет
малого ранга. CSSA работает как приближение: предполагается, что компоненте
сигнала соответствует одна из собственных троек траекторной матрицы с
максимальным сингулярным значением.

При типичной обработке одной прямой основная геометрическая структура сохраняется,
но после восстановления остаются слабые остаточные значения.

\begin{figure}[H]
	\centering
	\begin{subfigure}{0.49\textwidth}
		\centering
        \includegraphics[width=\textwidth]{singular-values-02.png}
		\caption{$\sigma=0.2$}
	\end{subfigure}
	\hfill
	\begin{subfigure}{0.49\textwidth}
		\centering
        \includegraphics[width=\textwidth]{singular-values-005.png}
		\caption{$\sigma=0.05$}
	\end{subfigure}
	\caption{Первые сингулярные значения для чистой и зашумленной строки после DFT}
	\label{fig:one-row-singular-values}
\end{figure}

Рис.~\ref{fig:one-row-singular-values} иллюстрирует, как шум нарушает точную
одноранговую структуру ряда. В чистом случае ненулевой остается только одна
компонента. При $\sigma=0.2$ шум растягивает сингулярный спектр, и
компонента сигнала может перестать быть первой. При меньшем шуме доминирование
первой компоненты выражено сильнее.

Наличие шума может приводить к смещению максимума: индекс наибольшего значения
в восстановленной строке не всегда совпадает с положением истинного пикселя.
Этот эффект связан с ошибкой выделения сигнальной компоненты на фоне шума, поэтому
при меньшем уровне шума он должен проявляться слабее. На рис.~\ref{fig:shift-profiles}
показаны две реконструкции для одного и того же положения сигнального пикселя:
при $\sigma=0.2$ максимум смещается, а при $\sigma=0.05$ основная масса
восстановленного сигнала концентрируется ближе к истинному индексу.

\begin{figure}[H]
	\centering
	\begin{subfigure}{0.49\textwidth}
		\centering
		\includegraphics[width=\textwidth]{shift_profiling.png}
		\caption{$\sigma=0.2$}
	\end{subfigure}
	\hfill
	\begin{subfigure}{0.49\textwidth}
		\centering
		\includegraphics[width=\textwidth]{shift_profiling_sigma_005.png}
		\caption{$\sigma=0.05$}
	\end{subfigure}
	\caption{График восстановленной строки при смещении максимума}
	\label{fig:shift-profiles}
\end{figure}

Другой характерный дефект --- растекание. Оно означает, что после обратного DFT
восстановленная строка имеет максимум в том же индексе, что и исходная, но
может содержать несколько соседних пикселей вместо одного с интенсивностью
$c_p$. Уменьшение уровня шума
частично ослабляет этот эффект, однако для дополнительного контроля растекания
можно использовать параметрическое восстановление на основе ESPRIT.
Пример такой реконструкции показан на рис.~\ref{fig:smear-profile}.

\begin{figure}[H]
	\centering
	\includegraphics[width=0.78\textwidth]{smear_profiling.png}
	\caption{График восстановленной строки при растекании сигнала}
	\label{fig:smear-profile}
\end{figure}

\subsection{ESPRIT как дополнение к \texorpdfstring{\texttt{CSSA ROW.ROW}}{CSSA ROW.ROW}}

Пусть $P_r$ --- это матрица, составленная из левых сингулярных векторов,
полученных в результате работы CSSA.
В идеальном случае, когда ряд является суммой $r$ комплексных экспонент,
его траекторная матрица имеет ранг $r$, а её столбцовое пространство
совпадает с пространством, порожденным векторами из комплексных экспонент:
\[
u(\mu_\ell)
=
\begin{pmatrix}
1 & \mu_\ell & \mu_\ell^2 & \cdots & \mu_\ell^{L-1}
\end{pmatrix}^{\top},
\qquad \ell=1,\dots,r,
\]
где $\mu_\ell=\exp(-2\pi\ii\omega_\ell)$. Иными словами,
\[
\operatorname{span}(P_r)
=
\operatorname{span}\{u(\mu_1),\dots,u(\mu_r)\}.
\]
Обозначим через $\underline P_r$ матрицу, полученную из $P_r$ удалением
последней строки, а через $\overline P_r$ --- матрицу, полученную из $P_r$
удалением первой строки.
Тогда
\[
\underline P_r,\overline P_r\in\mathbb C^{(L-1)\times r}.
\]
Если бы CSSA восстановил сигнал идеально, то полученная матрица имела бы вид
\[
P=
\begin{pmatrix}
1 & \cdots & 1\\
\mu_1 & \cdots & \mu_r\\
\vdots & \ddots & \vdots\\
\mu_1^{L-1} & \cdots & \mu_r^{L-1}
\end{pmatrix}\in\mathbb C^{L\times r},
\]
где
\[
\mu_j=\exp\!\left(-2\pi\ii\omega_j\right),\qquad j=1,\dots,r.
\]
и ее подматрицы $\underline P$ и $\overline P$, определенные аналогично
$\underline P_r$ и $\overline P_r$:
\[
\underline P=
\begin{pmatrix}
1 & \cdots & 1\\
\mu_1 & \cdots & \mu_r\\
\vdots & \ddots & \vdots\\
\mu_1^{L-2} & \cdots & \mu_r^{L-2}
\end{pmatrix},
\qquad
\overline P=
\begin{pmatrix}
\mu_1 & \cdots & \mu_r\\
\mu_1^2 & \cdots & \mu_r^2\\
\vdots & \ddots & \vdots\\
\mu_1^{L-1} & \cdots & \mu_r^{L-1}
\end{pmatrix}.
\]
Тогда выполняется точное соотношение
\[
\overline P=\underline P D,
\qquad
D=\operatorname{diag}(\mu_1,\dots,\mu_r).
\]
Для $P_r$ равенство, вообще говоря, нарушается из-за шума,
поэтому матрица перехода определяется как решение задачи минимизации нормы:
\[
\widehat M
=
\arg\min_{M\in\mathbb C^{r\times r}}
\|\overline P_r-\underline P_r M\|_F.
\]
Здесь $\|A\|_F$ --- норма Фробениуса матрицы $A$:
\[
\|A\|_F=\sqrt{\sum_{i,j}|a_{ij}|^2}.
\]
Если $\underline P_r$ имеет полный столбцовый ранг, то
\[
\widehat M=\underline P_r^\dagger\overline P_r,
\]
где $\underline P_r^\dagger$ обозначает псевдообратную матрицу
Мура--Пенроуза. В данном случае
$$\underline P_r^\dagger=(\underline P_r^{*}\underline P_r)^{-1}\underline P_r^{*}.$$
Далее используются собственные значения полученной матрицы $\widehat M$:
по ним восстанавливаются оценки модуля $\widehat\rho_i$ и частоты $\widehat\omega_i$ исходных экспонент.
Это приводит к следующему алгоритму.

\begin{algorithm}[H]
	\caption{ESPRIT}
	\begin{algorithmic}[1]
		\Require матрица $P_r\in\mathbb C^{L\times r}$
        \Ensure оценки параметров $(\widehat\rho_i,\widehat\omega_i)$, $i=1,\dots,r$
		\State построить матрицы $\underline P_r$ и $\overline P_r$
        \State найти матрицу $\widehat M\in\mathbb C^{r\times r}$ из соотношения
		\[
        \widehat M=\underline P_r^\dagger\overline P_r
		\]
        \State найти собственные значения $\widehat\mu_1,\dots,\widehat\mu_r$ матрицы $\widehat M$
		\State положить
		\[
        \widehat\rho_i=|\widehat\mu_i|,
		\qquad
        \widehat\omega_i=\operatorname{Arg}(\widehat\mu_i),
		\qquad i=1,\dots,r
		\]
        \State \Return $(\widehat\rho_i,\widehat\omega_i)$, $i=1,\dots,r$
	\end{algorithmic}
\end{algorithm}
После оценки параметров используется проекция оцененного на этапе CSSA сигнала на подпространство, порожденное
векторами
\[
u\bigl(\exp(-2\pi\ii\widehat\omega_\ell)\bigr),
\qquad \ell=1,\dots,r.
\]
Однако, не всякая комплексная экспонента при обратном DFT переходит в единичный орт.
\begin{statement}
Пусть $c\ne0, 0 \le \omega < 1$ и
\[
\widehat z_k
=
c\exp\!\left(-2\pi\ii\omega(k-1)\right),
\qquad
k=1,\dots, N_{\mathrm{col}}.
\]
Тогда
\[
\mathcal F^{-1}\widehat z=c\mathbf e_m
\]
для некоторого $m\in\{1,\dots,N_{\mathrm{col}}\}$ тогда и только тогда, когда:
\[
\omega = \frac{m-1}{N_{\mathrm{col}}}.
\]
\end{statement}

\begin{proof}
Если $\omega=(m-1)/N$, то
\[
\begin{aligned}
\bigl(\mathcal F^{-1}\widehat z\bigr)_j
&=
\frac{c}{N}
\sum_{k=1}^{N}
\exp\!\left(
2\pi\ii\left(\frac{j-1}{N_{\mathrm{col}}}-\omega\right)(k-1)
\right)
\\
&=
\frac{c}{N}
\sum_{k=1}^{N}
\exp\!\left(
2\pi\ii\frac{j-m}{N_{\mathrm{col}}}(k-1)
\right)
=
\begin{cases}
c, & j=m,\\
0, & j\ne m.
\end{cases}
\end{aligned}
\]
Значит, $\mathcal F^{-1}\widehat z=c\mathbf e_m$.

Обратно, пусть $\mathcal F^{-1}\widehat z=c\mathbf e_m$. Применяя DFT к обеим
частям, получаем
\[
c\exp\!\left(-2\pi\ii\omega(k-1)\right)
=
c\exp\!\left(
-2\pi\ii\frac{(m-1)(k-1)}{N_{\mathrm{col}}}
\right),
\qquad \forall k=1,\dots,N_{\mathrm{col}}.
\]
Так как $c\ne0$, при $k=2$ имеем
\[
\exp(-2\pi\ii\omega)
=
\exp\!\left(-2\pi\ii\frac{m-1}{N_{\mathrm{col}}}\right),
\]
то есть $\omega = (m-1)/N_{\mathrm{col}}$, в силу ограничений на $\omega$.
\end{proof}

Таким образом, для комплексной экспоненты растекание при обратном DFT
отсутствует тогда и только тогда, когда частота попадает в сетку шага
$1/N_{col}$. Как правило, это не так, поэтому после оценивания частоты было
предложено использовать ближайшее значение из элементов сетки:
$$
\widehat\omega_{\text{est}} = \frac{\left\lfloor \widehat\omega N_{\mathrm{col}} + 1/2 \right\rfloor}{N_{\mathrm{col}}}.
$$

\section{\texorpdfstring{\texttt{CSSA COL.ROW}}{CSSA COL.ROW} и ранг столбцов траекторной матрицы после DFT}

Для начала рассмотрим одну компоненту сигнала без шума. Ее матрицу
обозначим через $X$ и будем считать, что прямая задается целыми числами $a$ и $b$:
\begin{equation}
    X_{r,m}
    =
    A\mathbf 1_{\{r=am+b\}},
    \qquad
    r=1,\dots,N_{\mathrm{row}},
    \quad
    m=1,\dots,N_{\mathrm{col}}.
    \label{eq:X-definition}
\end{equation}
Здесь $a\in\mathbb N$, $b\in\mathbb Z$, $A\ne0$.

Обозначим через $\xi_{r,q}$ элементы матрицы после построчного DFT:
\begin{equation}
    \xi_{r,q}
    =
    \sum_{m=1}^{N_{\mathrm{col}}}
    X_{r,m}
    \exp\left(
        -2\pi\ii\frac{(m-1)(q-1)}{N_{\mathrm{col}}}
    \right),
    \qquad
    q=1,\dots,N_{\mathrm{col}}.
    \label{eq:dft-row}
\end{equation}
Зафиксируем индекс $q$ и обозначим соответствующий столбец через
$u$, то есть $u_r=\xi_{r,q}$, $r=1,\dots,N_{\mathrm{row}}$.

Именно к таким рядам $u$ применяется CSSA в варианте \texttt{COL.ROW}. Пусть
$K=N_{\mathrm{row}}-L+1$. Траекторной матрицей ряда $u$ будем называть
\begin{equation}
    \Hankel_L(u)
    =
    \{u_{p+\ell-1}\}_{p=1,\ell=1}^{L,K}
    =
    \begin{pmatrix}
        u_1 & u_2 & \cdots & u_K \\
        u_2 & u_3 & \cdots & u_{K+1} \\
        \vdots & \vdots & \ddots & \vdots \\
        u_L & u_{L+1} & \cdots & u_{N_{\mathrm{row}}}
    \end{pmatrix}
    \label{eq:trajectory-matrix}
\end{equation}

\begin{statement}
\label{st:colrow-main-rank}
Пусть $X$ задаётся формулой \eqref{eq:X-definition}, а $u$ --- фиксированный
столбец после построчного DFT \eqref{eq:dft-row}. Предположим, что выполнены условия
\begin{equation}
    a+\max(b,0)\le \min(L,K),
    \qquad
    K=N_{\mathrm{row}}-L+1,
    \label{eq:rank-fits-window}
\end{equation}
и условие:
\begin{equation}
    a(N_{\mathrm{col}}+1)+b>N_{\mathrm{row}}.
    \label{eq:no-bottom-truncation}
\end{equation}
Тогда
\begin{equation}
    \rank \Hankel_L(u)
    =
    a+\max(b,0).
    \label{eq:main-rank-formula}
\end{equation}
\end{statement}

\begin{proof}
Положим
\begin{equation*}
    \rho_q
    =
    \exp\left(
        -2\pi\ii\frac{q-1}{N_{\mathrm{col}}}
    \right).
\end{equation*}
Тогда
\begin{equation*}
    \exp\left(
        -2\pi\ii\frac{(m-1)(q-1)}{N_{\mathrm{col}}}
    \right)
    =
    \rho_q^{m-1}.
\end{equation*}
Поэтому из \eqref{eq:dft-row} получаем
\begin{equation*}
    u_r
    =
    \xi_{r,q}
    =
    \sum_{m=1}^{N_{\mathrm{col}}} X_{r,m}\rho_q^{m-1}.
\end{equation*}
Согласно \eqref{eq:X-definition}, матрица $X$ имеет ненулевые элементы только в тех строках, индексы которых имеют вид
\begin{equation*}
    r=am+b.
\end{equation*}
В этом случае
\begin{equation*}
    u_r=A\rho_q^{m-1}.
\end{equation*}
Следовательно, элементы столбца $u$ имеют вид
\begin{equation}
    u_r
    =
    \begin{cases}
        A\rho_q^{m-1}, & r=am+b,\quad m=1,\dots,N_{\mathrm{col}},\\
        0, & \text{иначе}.
    \end{cases}
    \label{eq:u-structure}
\end{equation}

Обозначим строки траекторной матрицы \eqref{eq:trajectory-matrix} через
\begin{equation*}
    R_1,\dots,R_L,
\end{equation*}
где
\begin{equation*}
    R_p=(u_p,u_{p+1},\dots,u_{p+K-1}).
\end{equation*}
Сначала рассмотрим случай $b\ge0$. Первая ненулевая строка ряда $u$
соответствует $m=1$. Положим
\begin{equation*}
    d=a+b.
\end{equation*}
Тогда эта строка имеет номер $d$. Так как $b\ge0$, то $d\ge a$. Первые $d-1$ строк ряда $u$ нулевые, а $d$-я строка ненулевая:
\begin{equation}
    u_1=\dots=u_{d-1}=0,
    \qquad
    u_d=A.
    \label{eq:first-nonzero}
\end{equation}
Поскольку $d-1=b+(a-1)$, первые $d-1$ нулей можно представить как $b$
начальных нулей и ещё $a-1$ нулей непосредственно перед первым ненулевым
элементом. Далее после каждого ненулевого элемента снова идут $a-1$ нулей,
так как соседние ненулевые строки отличаются на $a$:
\begin{equation*}
    u=
    (\underbrace{0,\dots,0}_{b},
    \underbrace{0,\dots,0}_{a-1},
    A,
    \underbrace{0,\dots,0}_{a-1},
    A\rho_q,
    \underbrace{0,\dots,0}_{a-1},
    A\rho_q^2,
    \dots).
\end{equation*}

Сначала докажем нижнюю оценку ранга. Рассмотрим левый верхний блок траекторной матрицы размера $d\times d$:
\begin{equation*}
    B
    =
    \{u_{p+\ell-1}\}_{p,\ell=1}^{d}
    =
    \begin{pmatrix}
        u_1 & u_2 & \cdots & u_d \\
        u_2 & u_3 & \cdots & u_{d+1} \\
        \vdots & \vdots & \ddots & \vdots \\
        u_d & u_{d+1} & \cdots & u_{2d-1}
    \end{pmatrix}.
\end{equation*}
Заметим, что из тождества (\ref{eq:first-nonzero}) следует, что $B$ имеет вид
\begin{equation*}
    B=
    \begin{pmatrix}
        0 & 0 & \cdots & 0 & A \\
        0 & 0 & \cdots & A & * \\
        0 & 0 & \ddots & * & * \\
        0 & A & \cdots & * & * \\
        A & * & \cdots & * & *
    \end{pmatrix},
\end{equation*}
где символ $*$ обозначает элементы, значения которых для данной оценки несущественны.

Переставим столбцы матрицы $B$ в обратном порядке. Получим матрицу
\begin{equation*}
    \widetilde B=
    \begin{pmatrix}
        A & 0 & 0 & \cdots & 0 \\
        * & A & 0 & \cdots & 0 \\
        * & * & A & \cdots & 0 \\
        \vdots & \vdots & \vdots & \ddots & \vdots \\
        * & * & * & \cdots & A
    \end{pmatrix}.
\end{equation*}
Эта матрица нижнетреугольная с ненулевой диагональю. Поэтому
\begin{equation*}
    \det \widetilde B=A^d\ne0.
\end{equation*}
Перестановка столбцов не меняет ранг, значит
\begin{equation*}
    \rank B=d.
\end{equation*}
Так как $B$ является подматрицей $\Hankel_L(u)$, получаем
\begin{equation*}
    \rank \Hankel_L(u)\ge d.
\end{equation*}

Теперь докажем обратное неравенство. Как видно из \eqref{eq:first-nonzero}, начиная с индекса $b+1$, ряд $u$ удовлетворяет соотношению
\begin{equation}
    u_{r+a}=\rho_q u_r,
    \qquad r\ge b+1.
    \label{eq:recursion-positive-b}
\end{equation}
Действительно, если $u_r\ne0$, то $r=am+b$ для некоторого $m\ge1$, и тогда
\begin{equation*}
    r+a=a(m+1)+b,
\end{equation*}
откуда
\begin{equation*}
    u_{r+a}=A\rho_q^m=\rho_q A\rho_q^{m-1}=\rho_q u_r.
\end{equation*}
Если же $u_r=0$, то $r \ne am+b$, то есть $r-b$ не делится на $a$. Поэтому $r+a-b$ также не делится на $a$, и
\begin{equation*}
    u_{r+a}=0=\rho_q u_r.
\end{equation*}

Используем этот факт, чтобы доказать, что строки матрицы $\Hankel_L(u)$ тоже
удовлетворяют рекуррентному соотношению. Рассмотрим произвольную строку $R_p$,
лежащую ниже первых $d$ строк траекторной матрицы, то есть пусть $p>d$. Так
как $d=a+b$, имеем
\begin{equation*}
    p-a>b,
\end{equation*}
а значит
\begin{equation*}
    p-a\ge b+1.
\end{equation*}
Для произвольного столбца $\ell=1,\dots,K$ элемент строки $R_{p-a}$ имеет вид $u_{p-a+\ell-1}$. Обозначим его индекс через
\begin{equation*}
    r=p-a+\ell-1.
\end{equation*}
Тогда $r\ge b+1$, поэтому к нему можно применить рекуррентное соотношение \eqref{eq:recursion-positive-b}:
\begin{equation*}
    u_{r+a}=\rho_q u_r.
\end{equation*}
Следовательно,
\begin{equation*}
    u_{p+\ell-1}
    =
    u_{r+a}
    =
    \rho_q u_r
    =
    \rho_q u_{p-a+\ell-1}.
\end{equation*}
Это верно для всех $\ell=1,\dots,K$, следовательно,
\begin{equation*}
    R_p=\rho_q R_{p-a},
    \qquad p>d.
\end{equation*}
Значит, каждая строка после $d$-й выражается через более раннюю строку.
Повторяя это соотношение, получаем, что все строки $R_{d+1},\dots,R_L$
принадлежат линейной оболочке первых $d$ строк.
Поэтому
\begin{equation*}
    \rank \Hankel_L(u)\le d.
\end{equation*}
Из нижней и верхней оценок следует
\begin{equation}
    \rank \Hankel_L(u)=d=a+b.
    \label{eq:rank-positive-b}
\end{equation}

Пусть теперь $b<0$. Первые значения $m$ могут давать строки $am+b<1$, то есть соответствующие точки находятся выше изображения и не попадают в матрицу $X$.

Пусть $m_0$ --- первый номер столбца, для которого точка попадает в изображение:
\begin{equation*}
    m_0=\min\{m\ge1:\ am+b\ge1\}.
\end{equation*}
Обозначим соответствующий индекс строки через
\begin{equation*}
    r_0=am_0+b.
\end{equation*}
Из определения $m_0$ имеем
\begin{equation*}
    r_0\ge1,
    \qquad
    r_0-a=a(m_0-1)+b<1,
\end{equation*}
поэтому
\begin{equation*}
    1\le r_0\le a.
\end{equation*}

Положим
\begin{equation*}
    C=A\rho_q^{m_0-1}\ne0.
\end{equation*}
Тогда ряд $u$ имеет вид
\begin{equation}
    u_r
    =
    \begin{cases}
        C\rho_q^t, & r=r_0+at,\quad t=0,1,2,\dots,\\
        0, & \text{иначе},
    \end{cases}
    \label{eq:u-negative-b-structure}
\end{equation}
Иными словами, после верхнего обрезания остаётся та же периодическая структура с шагом $a$, но первый ненулевой элемент находится в строке $r_0\in\{1,\dots,a\}$:
\begin{equation*}
    u=
    (\underbrace{0,\dots,0}_{r_0-1},
    C,
    \underbrace{0,\dots,0}_{a-1},
    C\rho_q,
    \underbrace{0,\dots,0}_{a-1},
    C\rho_q^2,
    \dots).
\end{equation*}

Из этой структуры следует рекуррентное соотношение
\begin{equation}
    u_{r+a}=\rho_q u_r.
    \label{eq:recursion-negative-b}
\end{equation}
для всех допустимых индексов. Поэтому для строк траекторной матрицы
\begin{equation*}
    R_{p+a}=\rho_q R_p,
    \qquad p=1,\dots,L-a.
\end{equation*}
Следовательно, все строки после первых $a$ выражаются через первые $a$ строк:
\begin{equation*}
    \rank \Hankel_L(u)\le a.
\end{equation*}

Теперь докажем нижнюю оценку. Рассмотрим левый верхний блок размера $a\times a$:
\begin{equation*}
    B
    =
    \{u_{p+\ell-1}\}_{p,\ell=1}^{a}
    =
    \begin{pmatrix}
        u_1 & u_2 & \cdots & u_a \\
        u_2 & u_3 & \cdots & u_{a+1} \\
        \vdots & \vdots & \ddots & \vdots \\
        u_a & u_{a+1} & \cdots & u_{2a-1}
    \end{pmatrix}.
\end{equation*}
Напомним, что в ряде $u$ первый ненулевой элемент находится в строке
$r_0$, где $1\le r_0\le a$, а расстояние между соседними ненулевыми
элементами равно $a$. Поэтому в любом фрагменте длины $a$ среди индексов
$1,\dots,2a-1$ находится ровно один ненулевой элемент ряда $u$.

Каждая строка блока $B$ как раз состоит из такого фрагмента:
\[
    u_p,\ u_{p+1},\dots,u_{p+a-1},
    \qquad p=1,\dots,a.
\]
Значит, в каждой строке блока $B$ стоит ровно один ненулевой элемент.
При переходе от строки $p$ к строке $p+1$ этот фрагмент сдвигается на один
индекс, поэтому положение ненулевого элемента в строке также сдвигается.
В результате ненулевые элементы попадают в разные столбцы, и в каждом
столбце блока $B$ тоже находится ровно один ненулевой элемент. Все эти
элементы имеют вид $C\rho_q^t$ и отличны от нуля, поэтому после перестановки
столбцов матрица $B$ становится диагональной с ненулевой диагональю. Значит,
\begin{equation*}
    \rank B=a.
\end{equation*}
Так как $B$ является подматрицей $\Hankel_L(u)$, получаем
\begin{equation*}
    \rank \Hankel_L(u)\ge a.
\end{equation*}
Из верхней и нижней оценок следует
\begin{equation}
    \rank \Hankel_L(u)=a.
    \label{eq:rank-negative-b}
\end{equation}

Объединяя \eqref{eq:rank-positive-b} и \eqref{eq:rank-negative-b}, получаем
\begin{equation}
    \rank \Hankel_L(u)
    =
    \begin{cases}
        a+b, & b\ge0,\\
        a, & b<0.
    \end{cases}
    \label{eq:piecewise-rank}
\end{equation}
Следовательно,
\begin{equation*}
    \rank \Hankel_L(u)
    =
    \max(a,a+b)
    =
    a+\max(b,0).
\end{equation*}
\end{proof}

Если условие
\(
    a(N_{\mathrm{col}}+1)+b>N_{\mathrm{row}}
\)
нарушается, в конце ряда $u$ появляются дополнительные элементы, и точная формула
\eqref{eq:main-rank-formula} уже не описывает общий случай. Поэтому далее
используется оценка ранга сверху.

\begin{statement}
\label{st:colrow-safe-rank-positive}
Пусть $X$ задаётся формулой \eqref{eq:X-definition} при $a\in\mathbb N$,
а $u$ --- фиксированный столбец после построчного DFT \eqref{eq:dft-row}.
Положим
\[
    y_1=a+b,
    \qquad
    y_N=aN_{\mathrm{col}}+b.
\]
Тогда
\begin{equation}
    \rank \Hankel_L(u)
    \le
    \widehat{R},
    \label{eq:safe-rank-bound}
\end{equation}
где
\begin{equation}
    \begin{aligned}
    \widehat{R}
    =
    \min\{&
        L,\ K,\
        \max(y_N,0),
        \max(N_{\mathrm{row}}+1-y_1,0),
        \\
        &
        a+\max(b,0)
        +
        \max\left(
            N_{\mathrm{row}}+1-a(N_{\mathrm{col}}+1)-b,\ 0
        \right)
    \}.
    \end{aligned}
    \label{eq:y-safe-rank}
\end{equation}
\end{statement}

\begin{proof}
Положим
\begin{equation*}
    \rho_q
    =
    \exp\left(
        -2\pi\ii\frac{q-1}{N_{\mathrm{col}}}
    \right).
\end{equation*}
Как и в доказательстве Утверждения~\ref{st:colrow-main-rank}, из
\eqref{eq:dft-row} получаем
\begin{equation}
    u_r
    =
    \begin{cases}
        A\rho_q^{m-1}, & r=am+b,\quad m=1,\dots,N_{\mathrm{col}},\\
        0, & \text{иначе}.
    \end{cases}
    \label{eq:u-structure-safe}
\end{equation}
\begin{enumerate}
    \item Траекторная матрица $\Hankel_L(u)$ имеет размер $L\times K$.
    Поэтому
    \begin{equation}
        \rank\Hankel_L(u)\le \min(L,K).
        \label{eq:rank-by-size}
    \end{equation}

    \item Поскольку $a\in\mathbb N$, индексы строк
    \[
        y_m=am+b,\qquad m=1,\dots,N_{\mathrm{col}},
    \]
    возрастают с ростом $m$. Поэтому $y_1=a+b$ является первым возможным
    индексом строки, в которой может стоять ненулевой элемент ряда $u$, а
    $y_N=aN_{\mathrm{col}}+b$ --- последним. Ненулевые элементы ряда $u$
    могут стоять только в строках
    \begin{equation*}
        y_1,\quad y_1+a,\quad \dots,\quad y_N.
    \end{equation*}
    С точностью до отсечения границами изображения ряд имеет вид
    \[
        u=(0,\dots,0,\ u_{y_1},\ 0,\dots,0,\ u_{y_1+a},\
        \dots,\ u_{y_N},\ 0,\dots,0),
    \]
    где индексы соседних ненулевых элементов отличаются на $a$.

    $p$-я строка траекторной матрицы состоит из элементов
    $u_p,u_{p+1},\dots,u_{p+K-1}$.
    Если $y_N<1$, то даже последний возможный ненулевой индекс меньше $1$,
    поэтому ряд $u$ нулевой. При $y_N\ge1$ из отсутствия ненулевых элементов
    после индекса $y_N$ следует схематический вид, где
    $s=\min(L,y_N)$:
    \[
        \Hankel_L(u)=
        \begin{pmatrix}
            u_1 & u_2 & \cdots & u_K\\
            \vdots & \vdots & \ddots & \vdots\\
            u_s & u_{s+1} & \cdots & u_{s+K-1}\\
            0 & 0 & \cdots & 0\\
            \vdots & \vdots & \ddots & \vdots\\
            0 & 0 & \cdots & 0
        \end{pmatrix},
    \]
    где нижний нулевой блок может быть пустым. Следовательно, ненулевых строк
    траекторной матрицы не больше $y_N$, и
    \begin{equation}
        \rank\Hankel_L(u)
        \le
        \max(y_N,0).
        \label{eq:rank-by-yN}
    \end{equation}

    Аналогично, если $y_1>N_{\mathrm{row}}$, то уже первый возможный индекс
    строки больше $N_{\mathrm{row}}$, поэтому ряд $u$ нулевой. Если
    $y_1\le N_{\mathrm{row}}$, то ненулевые элементы могут встречаться только
    при индексах $r\ge y_1$. Напомним, что
    \[
        \bigl(\Hankel_L(u)\bigr)_{p,\ell}=u_{p+\ell-1},
        \qquad
        p=1,\dots,L,\quad \ell=1,\dots,K.
    \]
    Значит $\ell$-й столбец траекторной матрицы состоит из элементов ряда с
    индексами
    \[
        \ell,\ \ell+1,\ \dots,\ \ell+L-1.
    \]
    Если $\ell+L-1<y_1$, то все эти индексы меньше $y_1$, поэтому весь
    $\ell$-й столбец нулевой. Иначе говоря, первые возможные ненулевые
    столбцы начинаются только с
    \[
        \ell_0=\max(1,\ y_1-L+1),
    \]
    и траекторная матрица имеет схематический вид
    \[
        \Hankel_L(u)=
        \begin{pmatrix}
            0 & \cdots & 0 & u_{\ell_0} & \cdots & u_K\\
            0 & \cdots & 0 & u_{\ell_0+1} & \cdots & u_{K+1}\\
            \vdots & \ddots & \vdots & \vdots & \ddots & \vdots\\
            0 & \cdots & 0 & u_{\ell_0+L-1} & \cdots & u_{N_{\mathrm{row}}}
        \end{pmatrix}.
    \]
    Здесь левый нулевой блок может быть пустым.
    Если $y_1>L$, то таких столбцов не больше
    \[
        K-(y_1-L+1)+1=N_{\mathrm{row}}+1-y_1.
    \]
    Если же $y_1\le L$, то потенциально ненулевых столбцов не больше
    $K=N_{\mathrm{row}}-L+1\le N_{\mathrm{row}}+1-y_1$. Следовательно, в
    обоих случаях ненулевых столбцов не больше $N_{\mathrm{row}}+1-y_1$, и
    \begin{equation}
        \rank\Hankel_L(u)
        \le
        \max(N_{\mathrm{row}}+1-y_1,0).
        \label{eq:rank-by-y1}
    \end{equation}

    \item Остаётся доказать оценку:
    \begin{equation}
        \rank\Hankel_L(u)
        \le
        a
        +
        \max(b,0)
        +
        \max\left(
            N_{\mathrm{row}}+1-a(N_{\mathrm{col}}+1)-b,\ 0
        \right).
        \label{eq:rank-by-ab}
    \end{equation}

    Индекс строки, который соответствовал бы следующему столбцу после
    $N_{\mathrm{col}}$, равен
    \begin{equation}
        y_{N+1}
        =
        a(N_{\mathrm{col}}+1)+b.
        \label{eq:y-next}
    \end{equation}
    Введём величину
    \begin{equation}
        h_+
        =
        \max\left(
            N_{\mathrm{row}}+1-y_{N+1},\ 0
        \right)
        =
        \max\left(
            N_{\mathrm{row}}+1-a(N_{\mathrm{col}}+1)-b,\ 0
        \right).
        \label{eq:h-plus}
    \end{equation}
    Интерпретация $h_+$ следующая: $y_{N+1}$ --- это первый индекс ряда $u$,
    в котором нарушается периодическая структура справа, тогда $h_+$ --- это
    количество нулей в конце ряда. Если $y_{N+1}>N_{\mathrm{row}}$, то
    $h_+=0$. Иначе все индексы
    $r\ge y_{N+1}$ занимают не более $h_+$ последних позиций ряда.

    Ряд $u$ устроен следующим образом: в начале стоят нули, связанные со
    сдвигом прямой, затем ненулевые элементы идут с шагом $a$, а после
    последней точки, соответствующей столбцам $1,\dots,N_{\mathrm{col}}$,
    может остаться нулевой хвост:
    \begin{equation*}
        u=
        (
        \underbrace{0,\dots,0}_{\max(b,0)},
        \quad
        \dots,\ *,
        \underbrace{0,\dots,0}_{a-1},
        *,
        \underbrace{0,\dots,0}_{a-1},
        *,\dots,
        \quad
        \underbrace{0,\dots,0}_{h_+}
        ).
    \end{equation*}
    Здесь $*$ обозначает значения вида $A\rho_q^{m-1}$. Эта структура
    нарушается на краях ряда: отсутствуют члены с $m\le0$ и с
    $m\ge N_{\mathrm{col}}+1$. Поэтому рассмотрим вспомогательный ряд с той
    же структурой без ограничения $1\le m\le N_{\mathrm{col}}$, а затем
    вычтем две краевые части.

    Введём вспомогательный конечный ряд $v=(v_1,\dots,v_{N_{\mathrm{row}}})$:
    \begin{equation}
        v_r
        =
        \begin{cases}
            A\rho_q^{m-1}, & r=am+b,\quad m\in\mathbb Z,\\
            0, & \text{иначе}.
        \end{cases}
        \qquad
        r=1,\dots,N_{\mathrm{row}}.
        \label{eq:v-definition}
    \end{equation}
    Для него выполнено рекуррентное соотношение для всех допустимых индексов:
    \begin{equation}
        v_{r+a}=\rho_q v_r.
        \label{eq:v-recursion}
    \end{equation}
    Следовательно, строки траекторной матрицы $\Hankel_L(v)$ удовлетворяют
    \begin{equation*}
        V_{p+a}=\rho_q V_p,
    \end{equation*}
    поэтому все строки выражаются через первые $a$ строк. Значит
    \begin{equation}
        \rank\Hankel_L(v)\le a.
        \label{eq:rank-v}
    \end{equation}

    Отделим в $v$ члены с номерами столбцов $m$ вне диапазона
    $1,\dots,N_{\mathrm{col}}$. Для $m\le0$ положим
    \begin{equation}
        w^-_r
        =
        \begin{cases}
            A\rho_q^{m-1}, & r=am+b,\quad m\le0,\\
            0, & \text{иначе}.
        \end{cases}
        \label{eq:w-minus}
    \end{equation}
    Если $b\le0$, то $w^-_r=0$ для всех $r$. Если $b>0$, то, в силу
    положительности $a$, наибольший
    возможный индекс строки у ненулевого элемента $w^-$ соответствует $m=0$
    и равен $b$. Поэтому траекторная матрица $\Hankel_L(w^-)$ имеет не более
    $\max(b,0)$ ненулевых строк, и
    \begin{equation}
        \rank\Hankel_L(w^-)
        \le
        \max(b,0).
        \label{eq:rank-w-minus}
    \end{equation}

    Для $m\ge N_{\mathrm{col}}+1$ положим
    \begin{equation}
        w^+_r
        =
        \begin{cases}
            A\rho_q^{m-1}, & r=am+b,\quad m\ge N_{\mathrm{col}}+1,\\
            0, & \text{иначе}.
        \end{cases}
        \label{eq:w-plus}
    \end{equation}
    Если $w^+_r\ne0$, то $r=am+b$ для некоторого
    $m\ge N_{\mathrm{col}}+1$. Минимальный такой индекс равен
    $a(N_{\mathrm{col}}+1)+b=y_{N+1}$. Значит внутри диапазона
    $1,\dots,N_{\mathrm{row}}$ ряд $w^+$ может быть ненулевым только не более
    чем на последних $h_+$ позициях. Следовательно,
    \begin{equation}
        \rank\Hankel_L(w^+)
        \le
        h_+
        =
        \max\left(
            N_{\mathrm{row}}+1-a(N_{\mathrm{col}}+1)-b,\ 0
        \right).
        \label{eq:rank-w-plus}
    \end{equation}

    По построению
    \begin{equation}
        u=v-w^- -w^+.
        \label{eq:u-v-decomposition}
    \end{equation}
    Траекторное отображение линейно, поэтому
    \begin{equation}
        \Hankel_L(u)
        =
        \Hankel_L(v)
        -
        \Hankel_L(w^-)
        -
        \Hankel_L(w^+).
        \label{eq:hankel-decomposition}
    \end{equation}
    Используя свойство ранга
    \[
        \rank(M_1+M_2+M_3)
        \le
        \rank M_1+\rank M_2+\rank M_3,
    \]
    и подставляя оценки \eqref{eq:rank-v}, \eqref{eq:rank-w-minus},
    \eqref{eq:rank-w-plus}, получаем \eqref{eq:rank-by-ab}.
\end{enumerate}

Из пунктов 1--3 следует \eqref{eq:safe-rank-bound}.

\end{proof}

От случая $a>0$ можно перейти к случаю $a<0$ с помощью отражения строк изображения относительно середины.

\begin{statement}
\label{st:colrow-safe-rank-integer}
Пусть в \eqref{eq:X-definition} допускается произвольное целое $a\ne0$.
Положим
\[
    \alpha=|a|.
\]
Если $a>0$, положим
\[
    y_1=a+b,
    \qquad
    y_N=aN_{\mathrm{col}}+b.
\]
Если $a<0$, положим
\[
    y_1=N_{\mathrm{row}}+1-(a+b),
    \qquad
    y_N=N_{\mathrm{row}}+1-(aN_{\mathrm{col}}+b).
\]
Тогда для фиксированного столбца $u$ после построчного DFT выполнена оценка
\begin{equation}
    \rank \Hankel_L(u)
    \le
    \widehat{R},
    \label{eq:integer-safe-rank-bound}
\end{equation}
где
\begin{equation}
    \begin{aligned}
    \widehat{R}
    =
    \min\{&
        L,\ K,
        \max(y_N,0),
        \max(N_{\mathrm{row}}+1-y_1,0),
        \\
        &
        \max(\alpha,y_1)
        +
        \max(N_{\mathrm{row}}+1-y_N-\alpha,0)
    \}.
    \end{aligned}
    \label{eq:integer-y-safe-rank}
\end{equation}
\end{statement}

\begin{proof}
При $a>0$ правая часть \eqref{eq:integer-y-safe-rank} совпадает с величиной
$\widehat{R}$ из Утверждения~\ref{st:colrow-safe-rank-positive}. Пусть теперь $a<0$.
Отразим строки изображения относительно середины:
\[
    \widetilde X_{r,m}=X_{N_{\mathrm{row}}+1-r,m}.
\]
Тогда линия $r=am+b$ переходит в линию
\[
    r=\alpha m+\widetilde b,
    \qquad
    \widetilde b=N_{\mathrm{row}}+1-b,
\]
с натуральным наклоном $\alpha=-a$. Для отражённой линии первый и последний
индексы строк равны
\[
    \alpha+\widetilde b
    =
    N_{\mathrm{row}}+1-(a+b)=y_1,
\]
и
\[
    \alpha N_{\mathrm{col}}+\widetilde b
    =
    N_{\mathrm{row}}+1-(aN_{\mathrm{col}}+b)=y_N.
\]
Поэтому величина $\widehat{R}$ из Утверждения~\ref{st:colrow-safe-rank-positive}
для отражённого изображения равна $\widehat{R}$.
Построчное DFT выполняется по столбцовому индексу $m$, поэтому отражение строк
с ним коммутирует: каждый фиксированный столбец после DFT просто переворачивается
относительно середины ряда,
\[
    \widetilde u_r=u_{N_{\mathrm{row}}+1-r}.
\]
При таком перевороте траекторная матрица получается из исходной изменением
порядка строк и столбцов на обратный, поэтому её ранг не меняется. Применяя
Утверждение~\ref{st:colrow-safe-rank-positive} к отражённому изображению, получаем
\eqref{eq:integer-safe-rank-bound}.
\end{proof}

Таким образом, при отсутствии шума для одной прямой получена достаточная оценка
числа компонент, обеспечивающая точное восстановление соответствующего
столбца после DFT. Тот же подход можно использовать для нескольких прямых. Если
\[
    u=\sum_{\ell=1}^{s} u^{(\ell)},
\]
где $u^{(\ell)}$ --- вклад $\ell$-й прямой в фиксированный столбец после DFT, то в
силу линейности оператора вложения
\[
    \Hankel_L(u)
    =
    \sum_{\ell=1}^{s}\Hankel_L\bigl(u^{(\ell)}\bigr).
\]
Следовательно,
\[
    \rank\Hankel_L(u)
    \le
    \sum_{\ell=1}^{s}
    \rank\Hankel_L\bigl(u^{(\ell)}\bigr).
\]
Поэтому для нескольких прямых можно выбирать число сохраняемых компонент не
меньше суммы достаточных оценок ранга для отдельных прямых. Такая оценка не
является необходимой, но дает корректное достаточное условие для применения
\texttt{CSSA COL.ROW} к сумме нескольких прямых.

\section{Пороговая фильтрация после CSSA}

Напомним, что после реконструкции после CSSA изображение обычно содержит не только
восстановленную линию, но и слабые остаточные значения. Для преобразования Хафа важна не
величина MSE, а активное множество
\[
E_{\tau}(\widehat{\mathbf{S}})
=
\{(i,j):\widehat S_{ij}>\tau\}.
\]
Поэтому после CSSA полезно вводить дополнительную пороговую фильтрацию. Один
из способов выбрать порог --- это оценить уровень шума в восстановленном
изображении. Для этого используем некоторые утверждения про шум и дискретное
преобразование Фурье.

\begin{statement}
Пусть $\boldsymbol{\varepsilon}\sim\mathcal{N}(0,\sigma^2 I_N)$ --- вещественный
гауссовский белый шум, а $F$ --- ненормированная матрица DFT:
\[
F_{kl}
=
\exp\!\left(
-\frac{2\pi\mathrm{i}}{N}(k-1)(l-1)
\right).
\]
Тогда $\widetilde{\boldsymbol{\varepsilon}}=F\boldsymbol{\varepsilon}$ является
комплексным гауссовским вектором,
\[
\mathbb{E}\widetilde{\boldsymbol{\varepsilon}}=0,
\qquad
\mathbb{E}\bigl(
\widetilde{\boldsymbol{\varepsilon}}
\widetilde{\boldsymbol{\varepsilon}}^{*}
\bigr)
=
N\sigma^2 I_N.
\]
\end{statement}

\begin{proof}
Линейное преобразование гауссовского вектора снова является гауссовским
вектором, поэтому $F\boldsymbol{\varepsilon}$ сохраняет гауссовскую природу.
Действительно, для любого комплексного вектора $a\in\mathbb C^N$ случайная
величина
\[
a^{*}F\boldsymbol{\varepsilon}
=
(F^{*}a)^{*}\boldsymbol{\varepsilon}
=
\operatorname{Re}\bigl((F^{*}a)^{*}\bigr)\boldsymbol{\varepsilon}
-
\ii\operatorname{Im}\bigl((F^{*}a)^{*}\bigr)\boldsymbol{\varepsilon}
\]
является комплексной случайной величиной, вещественная и мнимая части которой
являются линейными комбинациями координат вещественного гауссовского вектора
$\boldsymbol{\varepsilon}$. Значит, $a^{*}F\boldsymbol{\varepsilon}$ имеет
комплексное гауссовское распределение для любого $a$, то есть
$F\boldsymbol{\varepsilon}$ является комплексным гауссовским вектором.
Среднее равно нулю. Ковариационная матрица имеет вид
\[
\mathbb{E}\bigl(
F\boldsymbol{\varepsilon}
(F\boldsymbol{\varepsilon})^{*}
\bigr)
=
F\,\mathbb{E}(\boldsymbol{\varepsilon}\boldsymbol{\varepsilon}^{\top})\,F^{*}
=
\sigma^2 FF^{*}.
\]
При этом элемент произведения $FF^*$ с индексами $(k,\ell)$ равен
\[
\begin{aligned}
(FF^{*})_{k\ell}
&=
\sum_{s=1}^{N}
F_{ks}(F^{*})_{s\ell}
\\
&=
\sum_{s=1}^{N}
\exp\!\left(
-\frac{2\pi\ii}{N}(k-1)(s-1)
\right)
\exp\!\left(
\frac{2\pi\ii}{N}(\ell-1)(s-1)
\right)
\\
&=
\sum_{s=1}^{N}
\exp\!\left(
-\frac{2\pi\ii}{N}(k-\ell)(s-1)
\right)
=
\begin{cases}
N, & k=\ell,\\
0, & k\ne\ell,
\end{cases}
\end{aligned}
\]
поскольку при $k\ne\ell$ суммируется геометрическая прогрессия с множителем,
отличным от единицы. Значит, $FF^{*}=NI_N$, откуда и следует заявленная формула.
\end{proof}

\begin{corollary}
Пусть $\boldsymbol{\eta}$ --- комплексный гауссовский шум:
\[
\mathbb{E}\boldsymbol{\eta}=0,
\qquad
\mathbb{E}(\boldsymbol{\eta}\boldsymbol{\eta}^{*})
=
\widehat{\sigma}^{\,2}I_N.
\]
Тогда для ненормированного DFT
\[
\mathbb{E}\bigl(
\mathcal{F}^{-1}(\boldsymbol{\eta})
\mathcal{F}^{-1}(\boldsymbol{\eta})^{*}
\bigr)
=
\frac{\widehat{\sigma}^{\,2}}{N}I_N.
\]
\end{corollary}

\begin{proof}
Так как $\mathcal{F}^{-1}=N^{-1}F^{*}$, получаем
\[
\mathbb{E}\bigl(
\mathcal{F}^{-1}(\boldsymbol{\eta})
\mathcal{F}^{-1}(\boldsymbol{\eta})^{*}
\bigr)
=
\frac{1}{N^2}
F^{*}\,
\mathbb{E}(\boldsymbol{\eta}\boldsymbol{\eta}^{*})\,
F
=
\frac{\widehat{\sigma}^{\,2}}{N^2}F^{*}F
=
\frac{\widehat{\sigma}^{\,2}}{N}I_N.
\]
\end{proof}

Таким образом, один из вариантов состоит в том, чтобы использовать остаток для
задания порога. В SSA синусоидальная компонента является сигналом конечного
ранга, а белый шум при росте длины ряда асимптотически отделим от такого
сигнала \cite{ts_structure}. Поэтому после
выделения компонент CSSA, соответствующих сигналу, остаток естественно
рассматривать как оценку шума и использовать его для формирования порога $\tau$.
Пусть после DFT по строкам изображение представлено в виде
\[
\mathbf{Z}=\mathbf{S}+\boldsymbol{\eta},
\]
а после CSSA получена оценка $\widehat{\mathbf{S}}$. Тогда
\[
\mathbf{r}=\mathbf{Z}-\widehat{\mathbf{S}}
\]
является оценкой шумовой части. После обратного DFT положим
\[
\widehat D
=
\frac{1}{N}
\left\|
\mathcal{F}^{-1}(\mathbf{r})
\right\|_2^2,
\qquad
\widehat\sigma=\sqrt{\widehat D}.
\]
Тогда порог задается в виде
\[
\tau=c\,\widehat\sigma
=
c\sqrt{\widehat D}.
\]
Коэффициент $c$ подбирается экспериментально и зависит от уровня шума,
интенсивности линий и числа прямых.

На рис.~\ref{fig:threshold-row-col} показан пример действия такой фильтрации для
методов \texttt{CSSA ROW.ROW} и \texttt{CSSA COL.ROW}. Верхний ряд соответствует
восстановленным матрицам до пороговой обработки, нижний ряд --- матрицам после
пороговой фильтрации. После отсечения слабых остаточных значений активное множество
становится более разреженным, что уменьшает число посторонних точек, передаваемых
в преобразование Хафа.

\begin{figure}[H]
	\centering
	\begin{subfigure}{0.48\textwidth}
		\centering
		\includegraphics[width=\textwidth]{row_row_before_threshold.png}
		\caption{\texttt{ROW.ROW}, до порога}
	\end{subfigure}
	\hfill
	\begin{subfigure}{0.48\textwidth}
		\centering
		\includegraphics[width=\textwidth]{col_row_before_threshold.png}
		\caption{\texttt{COL.ROW}, до порога}
	\end{subfigure}

	\vspace{0.4em}

	\begin{subfigure}{0.48\textwidth}
		\centering
		\includegraphics[width=\textwidth]{row_row_after_threshold.png}
		\caption{\texttt{ROW.ROW}, после порога}
	\end{subfigure}
	\hfill
	\begin{subfigure}{0.48\textwidth}
		\centering
		\includegraphics[width=\textwidth]{col_row_after_threshold.png}
		\caption{\texttt{COL.ROW}, после порога}
	\end{subfigure}
	\caption{Пороговая фильтрация восстановленных матриц для \texttt{CSSA ROW.ROW} и \texttt{CSSA COL.ROW}}
	\label{fig:threshold-row-col}
\end{figure}

Другой вариант состоит в анализе гистограммы восстановленных значений. Если
на изображении несколько прямых с разными интенсивностями, то нужно сохранять
кластеры, соответствующие линиям, и удалять промежуточные остаточные значения.
Особенно важно не потерять слабую линию: если выделены несколько кластеров
линий по интенсивности, порог должен быть ниже среднего значения минимального
кластера, который соответствует линии малой интенсивности.

